% !TeX root = ../Book1.tex
\begin{figure}[htbp]
	\centering
	\begin{tikzpicture}[scale=1,font=\small]
		\begin{scope}[xshift=-3.5cm]
			\fill[blue!8] (-1.15,-1.05) rectangle (1.15,1.15);
			\draw[gray,->] (-1.2,0)--(1.25,0) node[right] {\(v^1\)};
			\draw[gray,->] (0,-1.1)--(0,1.3) node[above] {\(v^2\)};
			\draw[blue!70!black,thick,->] (0,0)--(0.7,0.7);
			\draw[blue!70!black,thick,->] (0,0)--(-0.65,-0.7);
			\fill (0,0) circle (1.4pt);
			\node at (0,-1.6) {\(\mathrm{T}_pM\)：全部方向};
		\end{scope}
		\begin{scope}
			\draw[gray,->] (-1.2,0)--(1.25,0) node[right] {\(v^1\)};
			\draw[gray,->] (0,-1.1)--(0,1.3) node[above] {\(v^2\)};
			\draw[blue!70!black,line width=1.2pt,<->] (-1,0)--(1,0);
			\fill (0,0) circle (1.4pt);
			\node at (0,-1.6) {双侧速度：\(v^2=0\)};
		\end{scope}
		\begin{scope}[xshift=3.5cm]
			\fill[blue!8] (-1.15,0) rectangle (1.15,1.15);
			\draw[gray,->] (-1.2,0)--(1.25,0) node[right] {\(v^1\)};
			\draw[gray,->] (0,-1.1)--(0,1.3) node[above] {\(v^2\)};
			\draw[blue!70!black,thick,->] (0,0)--(0.7,0.7);
			\draw[blue!70!black,thick,->] (0,0)--(-0.8,0.45);
			\draw[blue!70!black,line width=1.2pt,<->] (-1,0)--(1,0);
			\fill (0,0) circle (1.4pt);
			\node at (0,-1.6) {单侧速度：\(v^2\ge0\)};
		\end{scope}
	\end{tikzpicture}
	\caption[边界点的切空间与曲线速度]{二维边界点处的三种集合。左图是由导子定义的完整切空间；中图是双侧曲线的速度子空间；右图是单侧曲线的内向切锥。三图画的都是切向量分量，原点代表零向量。}
	\label{g1:fig:boundary-velocities}
\end{figure}
